% Discussion and Conclusion (draft 14 Nov 2026). Every statement is backed by docs/results_ledger.md; no citations needed here.

\section{Discussion and limitations}
\label{sec:discussion}
Sea surface temperature and sea level explain about half of the climatological error of the daily subsurface temperature (skill $+0.52$ and $+0.54$ on Test~A and Test~B), but most of that is already available to a per-cell linear regression. The U-Net ensemble adds a further $2$--$3\%$ in RMSE, almost all of it in the thermocline at 50--200~m, and adds nothing near the surface or below 300~m. We therefore regard the main result as a measurement of how much a neural network adds beyond a linear baseline under a protocol that cannot be tuned on the test years, not as an accuracy record.

The more consequential finding concerns the reference. The estimates are trained to reproduce GLORYS, and GLORYS is warmer than Argo in the thermocline in every year we examined, by an amount that changes over time. A model that agrees well with GLORYS therefore inherits this offset, and two designs rank in opposite order depending on the reference: the demeaned ensemble is the worst against GLORYS and the best against Argo. The yearly bias of the demeaned ensemble changes sign and its pattern error equals that of the frozen ensemble, so we read its advantage as bias reduction and not as a better estimate of the ocean. Evaluations against a reanalysis alone can hide such offsets.

\emph{Limitations.} (i)~The target is a reanalysis that assimilates Argo, so the Argo comparison is not independent and we make no claim of validation against two independent datasets. (ii)~The confidence intervals cover day-sampling variability only; they do not include training randomness, and the two test periods are single realisations. (iii)~Argo sampling is thin in the Andaman Sea (8 profiles on Test~B) and in 2025, and the Test~B target is the interim GLORYS production. (iv)~The demeaned ensemble was motivated by a sensitivity analysis on the validation window; it was pre-registered as co-primary before any test period was opened, and both ensembles are reported. (v)~The validation window was inspected 25 times before the models were frozen, the Test~A scoring was run twice (the second time only to add an output) and regional scoring was a further read of the test targets for added output; no model, input or recipe changed after a test period was opened. (vi)~Results are for one domain and one reanalysis.

\section{Conclusion}
\label{sec:conclusion}
Under a frozen chronological protocol, U-Net ensembles driven by SST and sea level anomalies reproduce GLORYS12 temperature in the North Indian Ocean with a pooled RMSE of 0.671$^\circ$C (Test~A) and 0.729$^\circ$C (Test~B), $1.8\%$ and $2.4\%$ below a per-cell linear regression, with the gain confined to the thermocline. Compared with Argo profiles the estimates agree to 0.834 and 0.909$^\circ$C, and the choice of reference changes which design looks best because the reanalysis itself drifts relative to Argo. Estimates trained on a reanalysis should therefore be checked against profiles, and gains over a linear baseline should be reported with intervals. Natural next steps, not attempted here, are to fine-tune or correct against Argo profiles directly and to test the approach in other basins.
